% Source: Wald GR, physical PDF pp. 31--36
%         (printed pp. 18--23).
% Coverage: Section 2.3, Tensors; the Metric Tensor; equations (2.3.1)--(2.3.10).
% Figures/tables/footnotes: footnote 2; no figures or tables.
% Status: source-reviewed and compile-clean.
% Uncertainties: none.

\subsection{Tensors; the Metric Tensor}\label{sec:2.3}

Given the notion of displacement vectors, the notion of tensors arises when one
considers other quantities of interest.  It turns out that many quantities have a
linear (or multilinear) dependence on displacements.  Consider, for example, a
measurement of the magnetic field (say, in the context of prerelativity physics).
For each probe orientation, a number is recorded: the magnetic field strength in
that given direction.  Since there is an infinite number of possible orientations of
the probe, in principle an infinite number of readings would be needed to determine
the magnetic field.  However, this is not necessary because the magnetic field
strength has a linear dependence on the probe orientation.  All that is required is
readings in three linearly independent probe directions; the reading in any other
probe direction is equal to a linear combination of these readings.

This fact gives rise to the notion of the magnetic field as a vector or, more
precisely, a covector.  One could define a covector as a collection of three numbers
(i.e., the probe readings) associated with a basis of spatial displacement vectors
(the three independent probe directions) which transform in an appropriate manner
when the basis is changed.  However, we will give below a simpler and more direct
definition of a covector as a linear map from spatial displacement vectors into
numbers.  We have defined the magnetic field as a covector here, but, as we shall
see at the end of this section, because space has a metric defined on it, we can
naturally associate to any covector an ordinary (spatial displacement) vector.

In a similar manner, other quantities that occur in physics have a similar linear
dependence on spatial displacement vectors but may be functions of more than one
such vector.  For example, for an ordinary material body in equilibrium, consider
the plane with normal vector \(\mathbf{n}\) passing through a point \(p\) in the body.
At \(p\), one could measure the force per unit area, \(F\), in the \(\mathbf{l}\)-direction
exerted on the matter on one side of the plane by the matter on the other side.  One
finds that \(F\) depends linearly on the choice of \(\mathbf{n}\) and \(\mathbf{l}\).  Thus,
although there is an infinite number of possible choices of \(\mathbf{n}\) and \(\mathbf{l}\),
the value of \(F\) for any \(\mathbf{n}\) and \(\mathbf{l}\) can be calculated by knowing
\(3\times3\) numbers, namely the values \(F\) takes when \(\mathbf{n}\) and \(\mathbf{l}\)
point in basis directions.  This motivates the definition of a tensor which will be
given below: A tensor is a multilinear (i.e., linear in each variable) map from
vectors (or covectors) into numbers.  The tensor which maps the pair of vectors
\((\mathbf{n},\mathbf{l})\) into the value of \(F\) is known as the \emph{stress tensor} of
the material body at \(p\).

We give now the precise mathematical definition of tensors and discuss their
properties.  Let \(V\) be any finite-dimensional vector space over the real
numbers.  (The case of prime interest for us is the tangent space, \(V=T_pM\).)
Consider the collection, \(V^\vee\), of linear maps \(f:V\to\mathbb{R}\).  If one
defines addition and scalar multiplication of such linear maps in the obvious way,
one gets a natural vector space structure on \(V^\vee\).  We call \(V^\vee\) the
\emph{dual space} to \(V\), and elements of \(V^\vee\) are called \emph{covectors}.
If \(e_1,\ldots,e_n\) is a basis of \(V\), we can define elements
\(f^{1},\ldots,f^{n}\in V^\vee\) by
\begin{equation}
  f^{\mu}(e_\nu)=\delta^\mu{}_{\nu} ,
  \tag{2.3.1}\label{eq:2.3.1}
\end{equation}
where \(\delta^\mu{}_{\nu}=1\) if \(\mu=\nu\) and 0 otherwise.  (This defines
the action of \(f^{\mu}\) on the basis elements; its action on an arbitrary vector
\(v\in V\) is determined by this and linearity.)  It follows directly (problem 6)
that \(\{f^{\mu}\}\) is a basis of \(V^\vee\), called the \emph{dual basis} to the
basis \(\{e_\mu\}\) of \(V\).  In particular, this shows that
\(\dim V^\vee=\dim V\).  The correspondence \(e_\mu\longleftrightarrow f^{\mu}\)
gives rise to an isomorphism between \(V\) and \(V^\vee\), but this isomorphism
depends on the choice of basis \(\{e_\mu\}\), so there is no natural way of
identifying \(V^\vee\) with \(V\) (unless more structure is given on \(V\), such as a
preferred basis or, as described below, a metric).

We now can apply the above construction starting with the vector space \(V^\vee\),
thereby obtaining the \emph{double dual space} to \(V\), denoted \(V^{\vee\vee}\).  An
element of \(V^{\vee\vee}\) is a linear map from \(V^\vee\) into \(\mathbb{R}\).  However,
\(V^{\vee\vee}\) is naturally isomorphic to the original vector space \(V\).  To each
vector \(v\) in \(V\) we can associate the map in \(V^{\vee\vee}\) whose value on the
covector \(w^\vee\in V^\vee\) is just \(w^\vee(v)\).  In this way, we obtain a one-to-one
linear map of \(V\) into \(V^{\vee\vee}\) which must be onto since
\(\dim V=\dim V^{\vee\vee}\).  Thus, taking the double dual gives nothing new; we can
naturally identify \(V^{\vee\vee}\) with the original vector space \(V\).  This
identification will be assumed in the discussion below.

We now are ready to define the notion of a tensor.  Let \(V\) be a finite
dimensional vector space and let \(V^\vee\) denote its dual space.  A tensor, \(T\),
of type \((k,l)\) over \(V\) is a multilinear map
\[
  T:\underbrace{V^\vee\times\cdots\times V^\vee}_{k}
  \times\underbrace{V\times\cdots\times V}_{l}\longrightarrow\mathbb{R} .
\]
In other words, given \(k\) covectors and \(l\) ordinary vectors, \(T\) produces a
number, and it does so in such a manner that if we fix all but one of the vectors or
covectors, it is a linear map on the remaining variable.

Thus, according to the above definition, a tensor of type \((0,1)\) is precisely a
covector.  Similarly, a tensor of type \((1,0)\) is an element of \(V^{\vee\vee}\).
However, since we identify \(V^{\vee\vee}\) with \(V\), a tensor of type \((1,0)\) is
nothing more than an ordinary vector.  Because of the identification of \(V^{\vee\vee}\)
with \(V\), we may view tensors of higher type in many different (though, of
course, equivalent) ways.  For example, a tensor \(T\) of type \((1,1)\) is a
multilinear map from \(V^\vee\times V\to\mathbb{R}\).  Hence, if we fix \(v\in V\),
\(T(\mathord\cdot,v)\) is an element of \(V^{\vee\vee}\), which we identify with an
element of \(V\).  Thus, given a vector in \(V\), \(T\) produces another vector in
\(V\) in a linear fashion.  In other words, we can view a tensor of type \((1,1)\)
as a linear map from \(V\) into \(V\), and vice versa.  Similarly, we can view
\(T\) as a linear map from \(V^\vee\) into \(V^\vee\).

With the obvious rules for adding and scalar multiplying maps, the collection
\(\mathcal T(k,l)\) of all tensors of type \((k,l)\) has the structure of a vector
space.  Because of the multilinearity property, a tensor is uniquely specified by
giving its values on vectors in a basis \(\{e_\mu\}\) of \(V\) and its dual basis
\(\{f^{\nu}\}\) of \(V^\vee\).  Since there are \(n^{k+l}\) independent ways of
filling the slots of a tensor of type \((k,l)\) with such basis vectors (where
\(n=\dim V=\dim V^\vee\)), the dimension of the vector space \(\mathcal T(k,l)\)
is \(n^{k+l}\).

We now introduce two simple but important operations on tensors, which will be
used commonly in what follows.  The first is called \emph{contraction} with
respect to the \(i\)th (covector) and \(j\)th (vector) slots and is a map
\(C:\mathcal T(k,l)\to\mathcal T(k-1,l-1)\) defined as follows.  If \(T\) is a
tensor of type \((k,l)\), then
\begin{equation}
  CT=\sum_{\sigma=1}^n
  T(\ldots,f^{\sigma},\ldots;\ldots,e_\sigma,\ldots) ,
  \tag{2.3.2}\label{eq:2.3.2}
\end{equation}
where \(\{e_\sigma\}\) is a basis of \(V\), \(\{f^{\sigma}\}\) is its dual basis,
and these vectors are inserted into the \(i\)th and \(j\)th slots of \(T\).  [Note
that the contraction of a tensor of type \((1,1)\), viewed as a linear map from
\(V\) into \(V\), is just the trace of the map.]  The tensor \(CT\) thus obtained
is independent of the choice of basis \(\{e_\mu\}\), so the operation of
contraction is indeed well defined (see problem 6).

The second operation on tensors is the outer product.  Given a tensor \(T\) of type
\((k,l)\) and another tensor \(T'\) of type \((k',l')\), we can construct a new
tensor of type \((k+k',l+l')\) called the \emph{outer product} of \(T\) and \(T'\)
and denoted \(T\otimes T'\), by the following simple rule.  Given \((k+k')\)
covectors \(f^{1},\ldots,f^{k+k'}\) and \((l+l')\) vectors
\(w_1,\ldots,w_{l+l'}\), we define \(T\otimes T'\) acting on these vectors to be
the product of \(T(f^{1},\ldots,f^{k};w_1,\ldots,w_l)\) and
\(T'(f^{k+1},\ldots,f^{k+k'};w_{l+1},\ldots,w_{l+l'})\).  Thus, one way of
constructing tensors is to take outer products of vectors and covectors.  A tensor
which can be expressed as such an outer product is called \emph{simple}.\footnote{In
many references, a tensor of type \((0,l)\) which can be expressed as the totally
antisymmetric part of a simple tensor (see Eq.~\eqref{eq:2.4.4} below) also is
referred to as simple.} If \(\{e_\mu\}\) is a basis of \(V\) and
\(\{f^{\nu}\}\) is its dual basis, it is easy to show that the \(n^{k+l}\) simple
tensors
\(\{e_{\mu_1}\otimes\cdots\otimes e_{\mu_k}\otimes f^{\nu_1}\otimes\cdots\otimes
f^{\nu_l}\}\) yield a basis of \(\mathcal T(k,l)\).  Thus, every tensor \(T\) of
type \((k,l)\) can be expressed as a sum of simple tensors in this collection
\begin{equation}
  T=\sum_{\mu_1,\ldots,\mu_k,\nu_1,\ldots,\nu_l=1}^n
  T^{\mu_1\cdots\mu_k}{}_{\nu_1\cdots\nu_l}
  e_{\mu_1}\otimes\cdots\otimes e_{\mu_k}\otimes
  f^{\nu_1}\otimes\cdots\otimes f^{\nu_l} .
  \tag{2.3.3}\label{eq:2.3.3}
\end{equation}
The basis expansion coefficients, \(T^{\mu_1\cdots\mu_k}{}_{\nu_1\cdots\nu_l}\),
are called the \emph{components} of the tensor \(T\) with respect to the basis
\(\{e_\mu\}\).  Note that we follow the standard convention in the notation for
components of \(T\) using superscripts for labels \(\mu_i\) associated with vectors
and subscripts for labels \(\nu_j\) associated with covectors.

In terms of components, we have the following formulas for contraction and outer
product.  Suppose the tensor \(T\) has components
\(T^{\mu_1\cdots\mu_k}{}_{\nu_1\cdots\nu_l}\) as in Eq.~\eqref{eq:2.3.3}.  Then,
the contraction, \(CT\), of \(T\) has components given by
\begin{equation}
  (CT)^{\mu_1\cdots\mu_{k-1}}{}_{\nu_1\cdots\nu_{l-1}}
  =\sum_{\sigma=1}^n
  T^{\mu_1\cdots\sigma\cdots\mu_{k-1}}{}_{
    \nu_1\cdots\sigma\cdots\nu_{l-1}} .
  \tag{2.3.4}\label{eq:2.3.4}
\end{equation}
If \(T'\) has components
\(T'^{\mu_1\cdots\mu_{k'}}{}_{\nu_1\cdots\nu_{l'}}\), then the outer product
\(S=T\otimes T'\) has components given by
\begin{equation}
  S^{\mu_1\cdots\mu_{k+k'}}{}_{\nu_1\cdots\nu_{l+l'}}
  =T^{\mu_1\cdots\mu_k}{}_{\nu_1\cdots\nu_l}
  T'^{\mu_{k+1}\cdots\mu_{k+k'}}{}_{\nu_{l+1}\cdots\nu_{l+l'}} .
  \tag{2.3.5}\label{eq:2.3.5}
\end{equation}

The above discussion applies to an arbitrary finite-dimensional vector space \(V\).
Let us now consider the case of prime interest for us, where \(V\) is the tangent
space, \(T_pM\), at point \(p\) of a manifold \(M\).  In this case, \(T_p^\vee M\) is
commonly called the \emph{cotangent space} at \(p\) and elements of \(T_p^\vee M\) are
called \emph{covectors}.  We also commonly refer to vectors in \(T_pM\) as
\emph{contravariant vectors} and covectors in \(T_p^\vee M\) as \emph{covariant
vectors}.  As discussed in Section~\ref{sec:2.2}, given a coordinate system, we can
construct a coordinate basis \(\partial/\partial x^1,\ldots,\partial/\partial x^n\)
of \(T_pM\).  The associated dual basis of \(T_p^\vee M\) is usually denoted as
\(\dd x^1,\ldots,\dd x^n\).  [Thus, we stress that \(\dd x^\mu\) is merely the symbol for
the linear map defined by
\(\dd x^\mu(\partial/\partial x^\nu)=\delta^\mu{}_{\nu}\).]  If we change
coordinate systems, we already showed that the components \(v'^{\mu'}\) of a
vector \(v\) in the new basis are related to its components \(v^\mu\) in the old
basis by the vector transformation law,
\begin{equation}
  v'^{\mu'}=\sum_{\mu=1}^n
  v^\mu\frac{\partial x'^{\mu'}}{\partial x^\mu} .
  \tag{2.3.6}\label{eq:2.3.6}
\end{equation}
If \(\omega_\mu\) denotes the components of a covector \(\omega\) with respect to
the dual basis \(\{\dd x^\mu\}\), then from Eqs.~\eqref{eq:2.3.1} and
\eqref{eq:2.3.6} it follows that under a coordinate transformation its components
become
\begin{equation}
  \omega'_{\mu'}=\sum_{\mu=1}^n
  \omega_\mu\frac{\partial x^\mu}{\partial x'^{\mu'}} .
  \tag{2.3.7}\label{eq:2.3.7}
\end{equation}
In general, the components of a tensor \(T\) of type \((k,l)\) transform as
\begin{equation}
\begin{aligned}
  T'^{\mu'_1\cdots\mu'_k}{}_{\nu'_1\cdots\nu'_l}
  &={\sum}_{\mu_1,\ldots,\mu_k,\nu_1,\ldots,\nu_l=1}^n
  T^{\mu_1\cdots\mu_k}{}_{\nu_1\cdots\nu_l}
  \frac{\partial x'^{\mu'_1}}{\partial x^{\mu_1}}\cdots
  \frac{\partial x'^{\mu'_k}}{\partial x^{\mu_k}}
  \frac{\partial x^{\nu_1}}{\partial x'^{\nu'_1}}\cdots
  \frac{\partial x^{\nu_l}}{\partial x'^{\nu'_l}} .
\end{aligned}
\tag{2.3.8}\label{eq:2.3.8}
\end{equation}
Equation~\eqref{eq:2.3.8} is known as the \emph{tensor transformation law}.

In other treatments, Eq.~\eqref{eq:2.3.8} often is used as the defining property of
a tensor.  The definition we have given here has the advantage that it generally is
much easier to define a quantity as a tensor by displaying it as a multilinear map on
vectors and covectors than it is to display it as a collection of numbers associated
with a coordinate system which changes according to Eq.~\eqref{eq:2.3.8} when we
change coordinate systems.  In fact, as we shall illustrate throughout this book, it
is rarely worthwhile to introduce a basis and take components of a tensor at all, let
alone to worry about how these components change under a change of basis.

An assignment of a tensor over \(T_pM\) for each point \(p\) in the manifold \(M\)
is called a \emph{tensor field}.  The notions of smoothness of a function and of a
(contravariant) vector field \(v\) were already defined in Section~\ref{sec:2.2}.  A
covector field \(\omega\) is said to be smooth (\(C^\infty\)) if for each smooth
vector field \(v\), the function \(\omega(v)\) is smooth.  A tensor \(T\) of type
\((k,l)\) is said to be smooth if for all smooth covector fields
\(\omega^1,\ldots,\omega^k\) and smooth contravariant vector fields
\(v_1,\ldots,v_l\), \(T(\omega^1,\ldots,\omega^k;v_1,\ldots,v_l)\) is a smooth
function.  The notion of a tensor field being \(C^k\) is defined similarly.

We now introduce the notion of a metric.  Intuitively, a metric is supposed to tell
us the ``infinitesimal squared distance'' associated with an ``infinitesimal
displacement.''  As discussed above in Section~\ref{sec:2.2}, the intuitive notion
of an ``infinitesimal displacement'' is precisely captured by the concept of a
tangent vector.  Thus, since ``infinitesimal squared distance'' should be quadratic
in the displacement, a metric, \(g\), should be a linear map from
\(T_pM\times T_pM\) into numbers, i.e., a tensor of type \((0,2)\).  In addition
the metric is required to be symmetric and nondegenerate.  By symmetric, we mean
that for all \(v_1,v_2\in T_pM\) we have \(g(v_1,v_2)=g(v_2,v_1)\).  By
nondegenerate, we mean that the only case in which we have \(g(v,v_1)=0\) for all
\(v\in T_pM\) is the case \(v_1=0\).  Thus, a \emph{metric}, \(g\), on a manifold
\(M\) is a symmetric, nondegenerate tensor field of type \((0,2)\).  In other
words, a metric is a (not necessarily positive definite) inner product on the tangent
space at each point.

In a coordinate basis, we may expand a metric \(g\) in terms of its components
\(g_{\mu\nu}\) as
\begin{equation}
  g=\sum_{\mu,\nu}g_{\mu\nu}\,\dd x^\mu\otimes \dd x^\nu .
  \tag{2.3.9}\label{eq:2.3.9}
\end{equation}
Sometimes the notation \(\dd s^2\) is used in place of \(g\) to represent the metric
tensor, in which case we write Eq.~\eqref{eq:2.3.9} as
\begin{equation}
  \dd s^2=\sum_{\mu,\nu}g_{\mu\nu}\,\dd x^\mu \dd x^\nu ,
  \tag{2.3.10}\label{eq:2.3.10}
\end{equation}
where, following standard practice, we have omitted writing the outer product sign
between \(\dd x^\mu\) and \(\dd x^\nu\).  The notation of Eq.~\eqref{eq:2.3.10}
conveys the intuitive flavor of a metric as representing ``infinitesimal squared
distance.''

Given a metric \(g\), we always can find an \emph{orthonormal basis}
\(e_1,\ldots,e_n\) of the tangent space at each point \(p\), i.e., a basis such that
\(g(e_\mu,e_\nu)=0\) if \(\mu\ne\nu\) and \(g(e_\mu,e_\mu)=\pm1\) (see problem
7).  There are, of course, many other orthonormal bases at \(p\), but the number of
basis vectors with \(g(e_\mu,e_\mu)=+1\) and the number with
\(g(e_\mu,e_\mu)=-1\) are independent of choice of orthonormal basis (problem 7).
The number of \(+\) and \(-\) signs occurring is called the \emph{signature} of the
metric.  In ordinary differential geometry, one usually deals with \emph{positive
definite} metrics, i.e., metrics with signature \(+\,+\,\ldots\,+\).  On the other
hand, the metric of spacetime has a signature \(-\,+\,+\,+\).  Positive definite
metrics are called \emph{Riemannian}; metrics with signatures like those on
spacetime (one minus and the remainder plus) are called \emph{Lorentzian}.

As defined above, at each point \(p\in M\) a metric \(g\) is a tensor of type
\((0,2)\) over \(T_pM\), i.e., a multilinear map from
\(T_pM\times T_pM\to\mathbb{R}\).  However, we can also view \(g\) as a linear
map from \(T_pM\) into \(T_p^\vee M\) via \(v\mapsto g(\mathord\cdot,v)\).  Because
of the nondegeneracy of \(g\), this map is one-to-one and onto.  In particular the
inverse map exists.  Thus, we can use \(g\) to establish a one-to-one correspondence
between vectors and covectors.  Indeed, given a metric \(g\) we could use this
correspondence to entirely circumvent the necessity of introducing covectors.
Normally this is done and accounts for why the concept of a covector is not more
familiar to most physicists.  However, in general relativity we shall be solving for
the metric of spacetime; since the metric is not known from the start, it is essential
that we keep the distinction between vectors and covectors completely clear.
